\uuid{OKUy}
\titre{Théorème de Bézout}
\auteur{bodin ; croizat ; sacre}
\organisation{exo7}
\licence{CC BY-NC-SA 4.0}
\datecreate{2022-03-08}
\niveau{L1}
\module{Algèbre}
\chapitre{Arithmétique dans Z}
\sousChapitre{Nombres premiers, nombres premiers entre eux}
\difficulte{1}
\theme{Théorème de Bézout}
\exo7qcmid{hilisit.arithmetique.016}

\contenu{
\begin{question}
\qclassification{103.04}
    Soient $a,b,c$ des entiers tels que $a | bc$.
    Dans le lemme de Gauss, quelle est la condition pour pouvoir conclure que $a|c$ ?
    \begin{answers}
        \good{$\pgcd(a,b)=1$}
        \bad{$\pgcd(a,c)=1$}
        \bad{$\pgcd(b,c)=1$}
        \bad{$a$, $b$ et $c$ sont des nombres premiers.}
    \end{answers}
    \begin{explanations}
    Lemme de Gauss : si $a | bc$ et $\pgcd(a,b)=1$ alors $a|c$.
    \end{explanations}
\end{question}
}
